\section*{الفصل \ref{ch:g12:doppler-effect} --- مفعول دوبلر}

\begin{solution}{exo:g12:doppler-effect:1}
(أ) $f' = 700 \times 340/315 \approx \qty{756}{Hz}$.
(ب) $f' = 700 \times 340/365 \approx \qty{652}{Hz}$.
(ج) \qty{700}{Hz}: فالسائق والصفّارة يتحركان معًا، ولا تتغير المسافة بينهما أبدًا.
\end{solution}

\begin{solution}{exo:g12:doppler-effect:2}
\emph{\omterm{def:g12:sound-acoustics:pitch-loudness}{العلوّ}} يخبو لأن \omterm{def:g10:signals-and-waves:wave}{الموجة} تنتشر --- فهو يتوقف على المسافة.
و\emph{\omterm{def:g12:sound-acoustics:pitch-loudness}{الحدّة}} تهبط من الاقتراب إلى الابتعاد --- وذلك \omterm{def:g12:doppler-effect:doppler}{مفعول دوبلر}، وهو يتوقف
على سرعة تغيّر المسافة. أي أثران مستقلان.
\end{solution}

\begin{solution}{exo:g12:doppler-effect:3}
راصد متحرك: $f' = 700\,(1 + 6.0/340) \approx \qty{712}{Hz}$؛ وبالركوب مبتعدًا،
$700\,(1 - 6.0/340) \approx \qty{688}{Hz}$. فالصفّارة ساكنة في الهواء، ومن ثم لا يضطرب نمط
\omterm{def:g10:signals-and-waves:wave}{الموجة}؛ إنما يركض الراصد خلاله.
\end{solution}

\begin{solution}{exo:g12:doppler-effect:4}
أمامه: $\lambda' = 310/400 = \qty{0.775}{m}$؛ وخلفه: $370/400 = \qty{0.925}{m}$؛ وفي حالة
السكون: $340/400 = \qty{0.85}{m}$. \omterm{def:g12:mechanical-waves:wavelength}{فطول الموجة} القصير يبلغ من يقترب منه القطار،
والطويل من يتركه.
\end{solution}

\begin{solution}{exo:g12:doppler-effect:5}
$\Delta f/f \approx 5.0/340 \approx 1.5\%$: $\Delta f \approx \qty{6.5}{Hz}$، أي نحو
\qty{446}{Hz} --- أي ربع نصف نغمة. وسيلاحظه موسيقي بالكاد.
\end{solution}

\begin{solution}{exo:g12:doppler-effect:6}
يستقبل الجدار $f\,v/(v - v_b)$ (بوصف المنبع متحركًا)؛ ويسمع الخفاش، وهو راصد متحرك يقترب من
\omterm{rem:g10:signals-and-waves:honest}{الصدى}، $(1 + v_b/v)$ ضعف ذلك: أي $f' = f(v + v_b)/(v - v_b) = 50.0 \times
346/334 \approx \qty{51.8}{kHz}$، بانزياح قدره $+\qty{1.8}{kHz}$. وهي صيغة
الانزياح المزدوج في مسدس \omterm{met:g10:signals-and-waves:echo}{الرادار} مع تبادل دورَي المُصدِر والعاكس: فلا تهمّ إلا سرعة الاقتراب.
\end{solution}

\begin{solution}{exo:g12:doppler-effect:7}
$u = c\,\Delta f/(2f) = \num{3.00e8} \times 4000/(2 \times \num{24.15e9}) \approx
\qty{24.8}{m/s} \approx \qty{89}{km/h}$. وعند \qty{50}{km/h} ($\qty{13.9}{m/s}$):
$\Delta f = 2 \times 13.9 \times \num{24.15e9}/\num{3.00e8} \approx \qty{2.2}{kHz}$.
\end{solution}

\begin{solution}{exo:g12:doppler-effect:8}
$u = v\,\Delta f/(2f) = 1540 \times 2600/\num{1.0e7} \approx \qty{0.40}{m/s}$.
أما الجريان العمودي فسرعته الشعاعية معدومة، ومن ثم لا انزياح البتة.
\end{solution}

\begin{solution}{exo:g12:doppler-effect:9}
(أ) المنبع: $f' = 1000 \times 340/306 \approx \qty{1111}{Hz}$. (ب) الراصد:
$f' = 1000 \times 1.1 = \qty{1100}{Hz}$. وتتنبأ الحركة البطيئة بالقيمة \qty{1100}{Hz}: وهي دقيقة في
حالة الراصد، وأقلّ بمقدار \qty{11}{Hz} في حالة المنبع --- أي تصحيح الرتبة $x^2$، وهو $1\%$ من
$f$ من أجل $x = 0.1$.
\end{solution}

\begin{solution}{exo:g12:doppler-effect:10}
$f = 2 \times 465 \times 415/880 \approx \qty{439}{Hz}$؛
$v_s = 340 \times 50/880 \approx \qty{19.3}{m/s} \approx \qty{70}{km/h}$. \omterm{def:g10:signals-and-waves:frequency}{والتواتر} الحقيقي $f$
هو المتوسط \omterm{def:g12:sound-acoustics:harmonics}{التوافقي}: فالاقتراب يرفع \omterm{def:g10:signals-and-waves:frequency}{التواتر} أكثر مما يخفضه الابتعاد، ومن ثم
يقع $f$ دون المتوسط الحسابي (\qty{440}{Hz}).
\end{solution}

\begin{solution}{exo:g12:doppler-effect:11}
طول \omterm{def:g10:signals-and-waves:wave}{موجة} أطول: أي \omterm{def:g12:doppler-effect:redshift}{انزياح نحو الأحمر}، فالنجم يبتعد. $v_r = c\,\Delta\lambda/\lambda =
\num{3.00e8} \times 0.2/656.3 \approx \qty{9.1e4}{m/s} \approx \qty{91}{km/s}$.
\end{solution}

\begin{solution}{exo:g12:doppler-effect:12}
$\Delta\lambda = \lambda K/c = 500 \times 55/\num{3.00e8} \approx \qty{9.2e-5}{nm}$؛
$\Delta\lambda/\lambda = 55/\num{3.00e8} \approx \num{1.8e-7}$. \omterm{def:g10:signals-and-waves:signal}{فالإشارة} كلها جزآن
من عشرة ملايين: \omterm{def:g10:light-spectra:white}{ومطياف} ينزاح بأكثر من $10^{-7}$ يدفنها.
\end{solution}

\begin{solution}{exo:g12:doppler-effect:13}
$v_r = 0.024\,c \approx \qty{7.2e6}{m/s}$. وسرعات تتناسب مع المسافة، وهي نفسها في
كل اتجاه: أي أن كل المسافات تتمدد تمددًا \omterm{def:g11:electric-gravitational-fields:capacitor}{منتظمًا} --- فكون متمدد بلا
مركز مميّز. وعند $2.4\%$ تكون التصحيحات من رتبة $(v_r/c)^2 \approx
\num{6e-4}$: فهي ما تزال جديرة بالثقة.
\end{solution}

\begin{solution}{exo:g12:doppler-effect:14}
تنزاح الحافتان في جهتين متعاكستين، ومن ثم يكون الخط الوسطي للقرص \emph{معرَّضًا} لا
منزاحًا: بعرض كامل قدره $2\lambda v/c = 2 \times 656.3 \times \num{2.0e3}/\num{3.00e8}
\approx \qty{8.8e-3}{nm}$ (نحو $\pm\qty{4.4e-3}{nm}$).
\end{solution}

\begin{solution}{exo:g12:doppler-effect:15}
(أ) أمام السيارة تتباعد القمم $(v - v_s)T$؛ وتفرّ الشاحنة بسرعة $v_o$ و
تلاقيها بالسرعة النسبية $v - v_o$: ومنه $f' = (v - v_o)/\bigl((v - v_s)T\bigr) =
f(v - v_o)/(v - v_s)$. (ب) $f' = 700 \times 310/300 \approx \qty{723}{Hz}$.
(ج) يعطي $v_o = v_s$ النتيجة $f' = f$: فالتباعد ثابت، وانعدام تغيّر المسافة
يعني انعدام انزياح دوبلر، بالتعريف.
\end{solution}

\begin{solution}{pb:g12:doppler-effect:1}
\textbf{1.} تتوسع كل قمة من النقطة التي صدرت منها: فتتزاحم القمم أمام
البوق المتحرك وتتمدد خلفه، ومن ثم يلاقيها الهاتف المقترب أسرع
($f_1$) من الهاتف المبتعد ($f_2$).

\textbf{2.} $f_1 = f\,\dfrac{v}{v - u}$، \quad $f_2 = f\,\dfrac{v}{v + u}$.

\textbf{3.} $f_1 - f_2 = \frac{2fv\,u}{v^2 - u^2}$ و$f_1 + f_2 =
\frac{2fv^2}{v^2 - u^2}$؛ وبالقسمة، $u = v\,\frac{f_1 - f_2}{f_1 + f_2} =
340 \times 60/882 \approx \qty{23}{m/s} \approx \qty{83}{km/h}$.

\textbf{4.} $\dfrac{2f_1 f_2}{f_1 + f_2} = f$ (فيتلاشى العامل المشترك):
$f = 2 \times 471 \times 411/882 \approx \qty{439}{Hz}$ --- أي المتوسط \omterm{def:g12:sound-acoustics:harmonics}{التوافقي}، وهو دون
المتوسط \qty{441}{Hz} لأن الاقتراب يرفع $f$ أكثر مما يخفضه الابتعاد.

\textbf{5.} أمامه: $(v - u)/f = 316.9/439 \approx \qty{0.72}{m}$؛ وخلفه:
$363.1/439 \approx \qty{0.83}{m}$؛ وفي حالة السكون: $340/439 \approx \qty{0.77}{m}$.

\textbf{6.} مرةً بوصفها \emph{راصدًا} متحركًا يستقبل \omterm{def:g10:signals-and-waves:wave}{الموجة}، ومرةً بوصفها
\emph{منبعًا} متحركًا يعيد إصدار ما استقبله.

\textbf{7.} $f' = f\,(1 + u/c)/(1 - u/c) \approx f(1 + 2u/c)$، ومنه
$\Delta f = f' - f \approx \dfrac{2u}{c}\,f$.

\textbf{8.} $u = c\,\Delta f/(2f) = \num{3.00e8} \times 3550/(2 \times \num{24.125e9})
\approx \qty{22}{m/s}$.

\textbf{9.} $22.1 \times 3.6 \approx \qty{79}{km/h}$: أي دون حدّ \qty{80}{km/h}
بقليل.

\textbf{10.} $\Delta f = 2f/(3.6\,c) \approx \qty{45}{Hz}$ لكل \unit{km/h}. وضرب
\omterm{rem:g10:signals-and-waves:honest}{الصدى} \omterm{def:g10:signals-and-waves:wave}{بالموجة} المصدَرة لا يبقي إلا \emph{الفارق}: أي تواترًا سمعيًّا
يُعدّ بلا عناء، وإن كان جزءًا قدره $10^{-7}$ من الحامل.

\textbf{11.} أطوال \omterm{def:g10:signals-and-waves:wave}{موجة} خطوط طيفية معلومة، ليلةً بعد ليلة؛ ثم
$v_r = c\,\Delta\lambda/\lambda$.

\textbf{12.} $K = \qty{55}{m/s}$، $P = \qty{4.2}{days}$.

\textbf{13.} $\Delta\lambda = 550 \times 55/\num{3.00e8} \approx \qty{1.0e-4}{nm}$،
$\Delta\lambda/\lambda \approx \num{1.8e-7}$ --- أي جزءًا من خمسة ملايين، أي أصغر
$\num{4e5}$ مرة من انزياح القطار $7\%$.

\textbf{14.} يتجاذب النجم والكوكب بالقدر نفسه، ومن ثم يدور كلاهما حول مركز
كتلتهما المشترك: فلا يقدر النجم أن يبقى ساكنًا. والمنحني الجيبي للسرعة $v_r$ \omterm{def:g11:forces-and-motion:ucm}{حركة دائرية
منتظمة} منظورًا إليها من الحافة --- أي \omterm{def:g10:universal-gravitation:satellite}{مدار} دائري.

\textbf{15.} لا تغيّر المسافة بين النجم والأرض إلا المركّبة على خط النظر، ومن ثم
لا يتغير \omterm{def:g12:mechanical-waves:wavelength}{طول الموجة} إلا بها. وفي المواجهة يكون $v_r = 0$ في كل حين: فلا \omterm{def:g10:signals-and-waves:signal}{إشارة} البتة.

\textbf{16.} $P = \qty{3.63e5}{s}$: $r^3 = \num{6.67e-11} \times \num{2.0e30} \times
(\num{3.63e5})^2/(4\pi^2) \approx \num{4.4e29}$، ومنه $r \approx \qty{7.6e9}{m}$ ---
أي نحو $1/20$ من المسافة بين الأرض والشمس.

\textbf{17.} $V = 2\pi r/P = 2\pi \times \num{7.6e9}/\num{3.63e5} \approx
\qty{1.3e5}{m/s}$.

\textbf{18.} $m = M K/V = \num{2.0e30} \times 55/\num{1.3e5} \approx \qty{8.3e26}{kg}$.

\textbf{19.} نحو $0.4$ من كتل المشتري (أي $\approx 140$ أرضًا)، تدور على مسافة أقرب $8$
مرة من عطارد: أي عملاق غازي يُشوى في وجه نجمه --- أي ``مشتري ساخن''.

\textbf{20.} الميل يخفي جزءًا من الحركة: فالسرعة $K$ المقيسة ليست إلا الحصة
الشعاعية، ومن ثم تكون الكتلة \omterm{def:g11:lenses-and-eye:images}{الحقيقية} \emph{أكبر} --- والقيمة \qty{8.3e26}{kg} حدّ أدنى. وانزياح دوري
قدره جزء من خمسة ملايين سلّم للتوّ مدارَ كوكب لم يره أحد قط وسرعته وكتلته الدنيا.
\end{solution}
